\chapter{Storage}
\label{ch:storage}

\cnans{Microsoft::Xna::Framework::Storage} is small --- three headers, two real classes and
one exception type --- and easy to overlook next to Graphics or Audio. It is also, per CNA's
own README, one of the namespaces reported at full presence alongside Graphics, Audio, and
Input/Touch, and it is worth its own short chapter precisely because it is the namespace
responsible for something every real game eventually needs: saving and loading user data.

\section{The XNA 4.0 fake-async pattern, once more}

\cnaclass{StorageDevice}'s own class documentation states its governing convention directly:
it ``uses the XNA 4.0 fake-async pattern: \texttt{BeginXxx} completes synchronously and the
paired \texttt{EndXxx} extracts the result.'' This is the same synchronous-but-async-shaped
pattern this book has already met repeatedly --- \cnaclass{LeaderboardReader}'s
\texttt{Begin}/\texttt{End} pairs (Chapter~5), \cnaclass{NetworkSession}'s own async
operations (Chapter~6), \cnaclass{AvatarDescription::BeginGetFromGamer} (Chapter~7) --- and
\cnaclass{StorageDevice} is arguably where the pattern originated in spirit: real XNA's
storage API had to accommodate a genuinely slow, UI-driven device-selection flow on a real
Xbox 360 (a player physically choosing which memory unit or hard drive to save to), even
though CNA's own desktop-filesystem-backed implementation completes every one of these calls
instantly.

\section{StorageDevice: free space, connection state, and device selection}

\cnaclass{StorageDevice} exposes \texttt{FreeSpace} and \texttt{TotalSpace} (both in bytes,
reporting real values from the underlying filesystem) and \texttt{IsConnected}. Opening a
container follows the fake-async pattern:
\texttt{BeginOpenContainer(displayName, callback, state)}/\texttt{EndOpenContainer(result)};
\texttt{DeleteContainer(titleName)} removes an entire container's directory tree outright.
Four overloads of the static \texttt{BeginShowSelector}/\texttt{EndShowSelector} pair cover
every combination real XNA supports --- with or without a specific \cnaclass{PlayerIndex},
with or without an explicit required free-space and directory-count --- for showing the
device-selection UI, even though, again, there is only ever one real ``device'' (the local
filesystem) for this to resolve to.

Two \texttt{NOXNA} static methods exist specifically because CNA has no Xbox-360-style
concept of a physically distinct storage device to select in the first place:
\texttt{SetAppNameEXT(appName)} sets the application name used to build a real,
per-application storage root directory on the host filesystem, and
\texttt{GetStorageRootEXT()} returns that resolved root path. A game is expected to call
\texttt{SetAppNameEXT} once, early --- the class's own documentation recommends the
\cnaclass{Game} constructor --- before any other storage call, since the storage root is
lazily resolved and cached on first use.

\section{StorageContainer: a real, if modest, virtual filesystem}

\cnaclass{StorageContainer} inherits \texttt{System::Object} and \texttt{System::IDisposable},
and provides exactly the file and directory operations a save-data system needs, all relative
to the container's own root: \texttt{CreateDirectory}/\texttt{DirectoryExists}/%
\texttt{DeleteDirectory}, \texttt{GetDirectoryNames()} (with a glob-pattern overload),
\texttt{CreateFile}/\texttt{FileExists}/\texttt{DeleteFile}, \texttt{GetFileNames()} (with the
same glob-pattern overload), and a three-overload \texttt{OpenFile} family covering file mode
alone, mode plus access, and mode plus access plus sharing --- each returning a real
\cnaclass{System::IO::Stream} (Chapter~10 of this volume covers \texttt{sharp-runtime}'s own
\cnaclass{Stream} type) for a game to read or write through directly. Every relative path
passed to any of these methods is resolved against the container's own root, which is itself
resolved against \texttt{StorageDevice::GetStorageRootEXT()} --- meaning the entire
namespace, end to end, ultimately roots itself in that one \texttt{NOXNA} extension point.

\section{StorageDeviceNotConnectedException}

\cnaclass{StorageDeviceNotConnectedException} is the one exception type this namespace adds,
thrown when an operation requires a connected storage device that is not currently available
--- present for API completeness with real XNA's own exception surface, even though CNA's own
always-available local-filesystem backing means this condition is rare in practice compared
to a real Xbox 360, where a memory unit could genuinely be physically removed mid-session.

\section{Why this namespace is worth remembering}

\cnans{Storage} is a good, compact illustration of a pattern this book has now shown several
times: an XNA namespace originally designed around real, physical, removable hardware and a
console's own UI conventions, faithfully reproduced in API shape, but backed by a
straightforward, always-available local filesystem underneath. The fake-async pattern, the
device-selection UI stubs, and the \texttt{NOXNA} app-name/storage-root extension together
tell the same story Chapter~1 of Volume~I told about the project as a whole: preserve the
real API's shape precisely, even where the platform reality underneath it is simpler than
what the API was originally built to abstract over.
